% Part: incompleteness
% Chapter: theories-computability
% Section: q-is-ce

\documentclass[../../../include/open-logic-section]{subfiles}

\begin{document}

\olfileid{inc}{tcp}{qce}

\olsection{$\Th{Q}$، \printtoken{S}{c.e.}-بشپړه ده}


\begin{thm}
$\Th{Q}$، !!{c.e.} ده خو د پرېکړې وړ نۀ ده۔ په حقيقت کښې دا يو
بشپړ !!{c.e.} سټ دے۔
\end{thm}

\begin{proof}
دا ليدل ستونزمن نۀ دي چې~$\Th{Q}$، !!{c.e.} ده، ځکه دا د هغو
جملو~$y$ د (کوډونو) سټ دے چې داسې~$x$ وي چې د~$y$ په~$\Th{Q}$ کښې
ثبوت وي:
\[
Q = \Setabs{y}{\lexists[x][\Prf[\Th{Q}](x,y)]}.
\]
خو پوهېږو چې~$\Prf[\Th{Q}](x,y)$ محاسبه کېدونکې ده (په حقيقت کښې
بنسټيزه بازګشتي ده)، او هر هغه سټ چې په پورتنۍ بڼه وليکل شي، په
محاسبوي ډول د شمېر وړ وي۔

دا ويل چې دا يو بشپړ په محاسبوي ډول د شمېر وړ سټ دے، له دې خبرې
سره برابر دي چې~$K \leq_m Q$، چېرته چې
$K = \Setabs{x}{!A_x(x) \downarrow}$۔ نو راځئ وښيو چې~$K$،
$\Th{Q}$ ته راکمېدونکے دے۔ ځکه چې د کلېني پريديکات~$T(e,x,s)$
بنسټيزه بازګشتي دے، په~$\Th{Q}$ کښې، فرض کړئ، د~$!A_T$ په وسيله
تمثيلېږي۔ بيا د هر~$x$ دپاره لرو:
\begin{align*}
x \in K & \lif \lexists[s][T(x,x,s)] \\
& \lif \lexists[s][(\Th{Q} \Proves !A_T(\num x,\num x, \num s))]
  \\
& \lif \Th{Q} \Proves \lexists[s][!A_T(\num x, \num x, s)].
\end{align*}
برعکس، که~$\Th{Q} \Proves \lexists[s][!A_T(\num x, \num x, s)]$،
نو په حقيقت کښې د کوم طبيعي عدد~$n$ دپاره بايد فارمول
$!A_T(\num x, \num x, \num n)$ رښتيا وي۔ اوس که~$T(x,x,n)$ ناسم
وای، نو ځکه چې~$!A_T$، $T$ تمثيلوي، $\Th{Q}$ به
$\lnot !A_T(\num x, \num x, \num n)$ ثابت کړے وای۔ خو بيا به
$\Th{Q}$ يو ناسم فارمول ثابت کړے وای، چې تناقض دے۔ نو~$T(x,x,n)$
بايد رښتيا وي، او له دې څخه~$!A_x(x) \downarrow$ لازمېږي۔

په لنډ ډول، د هر~$x$ دپاره، $x$ هله او يوازې هله په~$K$ کښې دے
چې~$\Th{Q}$، $\lexists[s][!A_T(\num x,\num x,s)]$ ثابت کړي۔ نو هغه
تابعه~$f$ چې~$x$ د جملې~$\lexists[s][!A_T(\num x,
  \num x, s)]$ (يو کوډ) ته وړي، له~$K$ څخه~$\Th{Q}$ ته راکمونه ده۔
\textbf{د ژباړې د نسخې سمون (OLCMP-091):} سرچينې په دې وروستۍ
پايله کښې د تمثيلوونکي فارمول پر ځاے د ماوراءژبې اړيکه ليکلې وه؛
دلته هماغه فارمول کارول شوے چې په مخکښې استدلال کښې معرفي او ثابت
شوے دے۔
\end{proof}

\end{document}
